% Gauss, Werke X/1 (1917), native77–81 upper, printed70–74 upper.
\SourceLeaf{077}\NoSourceHead
\SourceGroup{XII}
\HistoricTitle{[ZUR GESCHICHTE DER THEORIE DER KUBISCHEN UND BIQUADRATISCHEN RESTE.]}
\begin{center}\large\textsc{Gauss} an \textsc{Sophie Germain}.\end{center}
\par\indent Votre lettre du 20 Fevrier, mais qui ne m’est parvenue que le 12 Mars,
a été pour moi la source d’autant de plaisir que de surprise. Combien l’acquisition
d’une amitié aussi flateuse et précieuse est-elle douce à mon coeur.
L’interêt vif, que Vous avez pris à mon sort pendant cette guerre funeste merite
la plus sincere reconnoissance. Assurément, Votre lettre au Général
\textsc{Pernety} m’eût été fort utile, si j’avois été dans le cas d’avoir recours à une
protection specielle de la part du gouvernement françois. Heureusement les
evenemens et les suites de la guerre ne m’ont pas touché de trop près jusqu’ici,
bienque je sois persuadé, qu’elles auront une grande influence sur le plan futur
de ma vie. Mais comment Vous décrire mon admiration et mon étonnement,
en voïant se métamorphoser mon correspondant estimé Mr. \textsc{Leblanc} en cet
illustre personnage, qui donne un exemple aussi brillant de ce que j’aurois
peine de croire. Le goût pour les sciences abstraites en general et surtout
pour les mysteres des nombres est fort rare: on ne s’en étonne pas; les charmes
enchanteurs de cette sublime science ne se decelent dans toute leur beauté
qu’à ceux qui ont le courage de l’approfondir. Mais lorsqu’une personne de
ce sexe, qui, par nos moeurs et par nos prejugés, doit rencontrer infiniment
plus d’obstacles et de difficultés, que les hommes, à se familiariser avec ces
recherches epineuses, sait neanmoins franchir ces entraves et penetrer ce
qu’elles ont de plus caché, il faut sans doute, qu’elle ait le plus noble courage,
des talens tout à fait extraordinaires, le génie superieur. En effet rien
ne pourroit me prouver d’une maniere plus flateuse et moins equivoque, que
\clearpage
\SourceLeaf{078}\ArithmeticHead{71}{R}{GAUSS AN SOPHIE GERMAIN.}
\noindent les attraits de cette science, qui a embelli ma vie de tant de jouissances, ne
sont pas chimeriques, que la predilection, dont Vous l’avez honorée.

Les notes savantes, dont toutes Vos lettres sont si richement remplies,
m’ont donné mille plaisirs. Je les ai étudiées avec attention, et j’admire la
facilité, avec laquelle Vous avez penetré toutes les branches de l’Arithmetique,
et la sagacité avec laquelle Vous les avez su generaliser et perfectionner. Je
Vous prie d’envisager comme une preuve de cette attention, si j’ose ajouter
une remarque à un endroit de Votre derniere lettre. Il me semble, que la
proposition inverse, savoir »\textls[100]{si la somme des puissances }$n^{\mathrm{emes}}$\textls[100]{ de deux
nombres quelconques est de la forme }$hh+nff$\textls[100]{, la somme de ces nombres
eux-mêmes sera de la meme forme}« est enoncée un peu trop generalement.
Voici un exemple où cette règle est en défaut:
\[
\begin{gathered}
15^{11}+8^{11}=8649755859375+8589934592={}\\
8658345793967=1595826^2+11\mathpunct{.}745391^2
\end{gathered}
\]
\noindent Neanmoins $15+8=23$ ne peut se reduire sous la forme $xx+11yy$.

Il en est de même de la proposition: \textls[100]{si l’un des facteurs de la formule }
$yy+nzz$ \textls[100]{(}$n$\textls[100]{ étant un nombre premier) est de la forme }$(1,0,n)$\textls[100]{,
l’autre appartient necessairement à la meme forme.} Votre demonstration
ne prouve que ce, qu’aucune autre forme \textls[100]{indéfinie}, que telle qui
est equivalente à $(1,0,n)$, multipliée par la forme $(1,0,n)$, ne peut donner le
produit $(1,0,n)$, mais cette demonstration ne s’étend pas sur les nombres
\textls[100]{définis}. Soit, pour le déterminant $-n$, $C$ une classe de formes, quelconque
mais ni equivalente à la principale, ni à une autre classe \textls[100]{anceps}, soit $D$
la classe resultante de la duplication de $C$ (qui sera differente de la principale),
enfin soit $D'$ la classe opposée à $D$. Il s’ensuit, que de la composition
de $C+C+D'$ résulte la classe principale. Ainsi si les deux nombres $f,g$
peuvent être representés par une forme de la classe $C$, et le nombre $h$ par
une forme de la classe $D'$, le produit $fg\times h$ peut se reduire à $(1,0,n)$; mais
il est facile [de voir] que $fg$ ne se reduit pas seulement à $D$ ou $D'$ mais aussi à
$(1,0,n)$. Nous avons donc ici le cas, qu’un facteur $fg$, et le produit $fg\mathpunct{.}h$ sont
de la forme $(1,0,n)$, sans que pourtant l’autre facteur y appartienne nécessairement.
Au reste on voit facilement que le premier facteur doit être composé,
\clearpage
\SourceLeaf{079}\ArithmeticHead{72}{L}{ARITHMETIK.\quad NACHLASS.}
\noindent sans cela la proposition serait juste. Dans l’exemple ci dessus le facteur
$\dfrac{15^{11}+8^{11}}{23}$ enveloppe le diviseur $67$.

Depuis cinq ans des travaux astronomiques — auxquels pour le dire en
passant je dois surtout l’heureuse situation, dont j’ai joui pendant la vie de
notre duc, le victime malheureux de son attachement fidel à la maison de
Prusse — m’ont empeché de me delivrer autant qu’auparavant à ma predilection
pour l’Arithmetique et les autres branches de l’analyse. Je n’ai pas
pourtant negligé celle ci tout à fait. Tout au contraire j’ai rassemblé peu à
peu un grand nombre de recherches, qui un jour fourniront un autre volume
— sinon deux — certainement pas moins interessant que le premier. Même
dans le dernier hiver j’ai reussi à y ajouter une branche entierement nouvelle.
C’est la theorie des residus cubiques et des residus biquarrés, portée à un
degré de perfection égal à celui, qu’a atteint la theorie des residus quarrés.
Je mets cette theorie, qui repand un nouveau jour sur les residus quarrés
parmi les recherches les plus curieuses donc je me sois jamais occupé. Je
ne saurais Vous en donner une idée sans écrire un Memoire expres. Voici
pourtant quelque theoreme special, qui pourra servir d’un petit echantillon.

I. Soit $p$ un nombre premier de la forme $3n+1$. Je dis, que $2$ (c. à. d.
$+2$ et $-2$) est residu cubique de $p$, si $p$ se reduit à la forme $xx+27yy$;
que $2$ est Non-residu cubique de $p$, si $4p$ se reduit à cette forme. P. E. 7.
13. 19. 31. 37. 43. 61. 67. 73. 79. 94. Vous ne trouveres que $31=4+27$,
$43=16+27$, et $2\equiv4^3\ (\mathrm{mod.}\ 31)$, $2\equiv(-9)^3\ (\mathrm{mod.}\ 43)$.

II. Soit $p$ un nombre premier de la forme $8n+1$. Je dis que $+2$ et
$-2$ seront residus ou non-residus biquarrés de $p$, suivant ce que $p$ est ou
n’est pas de la forme $xx+64yy$. Par ex. parmi les nombres 17. 41. 73. 89.
97. 113. 137 Vous ne trouves que $73=9+64$, $89=25+64$, $113=49+64$,
et $25^4\equiv2\ (\mathrm{mod.}\ 73)$, $5^4\equiv2\ (\mathrm{mod.}\ 89)$, $20^4\equiv2\ (\mathrm{mod.}\ 113)$.

Les demonstrations de ces theoremes et de ceux qui sont plus generaux
sont intimement liées à des recherches delicates. — Voici une autre proposition
relative aux residus quarrés, dont la demonstration est moins cachée: je ne
l’ajoute pas, pour ne pas Vous derober le plaisir de la developper Vous-même,
si Vous la trouveres digne d’occuper quelques momens de Votre loisir.
\clearpage
\SourceLeaf{080}\ArithmeticHead{73}{R}{GAUSS AN SOPHIE GERMAIN.}
\par\indent Soit $p$ un nombre premier. Soient les $p-1$ nombres inferieurs à $p$ partagés
en deux classes
\[
\begin{aligned}
A\;&.\ .\ .\ .\ .\ .\ .\ .\ .\;1,2,3,4,\ldots,\tfrac12(p-1)\\
B\;&.\ .\ .\ .\ .\ .\ .\ .\ .\;\tfrac12(p+1),\ \tfrac12(p+3),\ \tfrac12(p+5),\ldots,p-1.
\end{aligned}
\]
\noindent Soit $a$ un nombre quelconque non-divisible par $p$. Multipliés tous les nombres
$A$ par $a$; prenés en les moindres residus selon le module $p$, soient,
entre ces residus, $\alpha$ appartenants à $A$ et $\beta$ appartenants à $B$ de sorte que
$\alpha+\beta=\tfrac12(p-1)$. Je dis, que $a$ est residu quarré de $p$ lorsque $\beta$ est pair,
non-residu lorsque $\beta$ est impair.

On peut tirer de cette proposition plusieurs consequences très remarquables;
entre-autres, elle donne le moïen d’etendre l’induction, par laquelle
on rassemble des cas speciels du theoreme fondamental \textls[100]{aussi loin qu’on
veut}, ce qui ne pourroit se faire par les methodes exposés art. 116—124[*)].

J’ai donné dans mon ouvrage deux demonstrations rigoureuses de ce fameux
theoreme et j’en possede encore \textls[100]{trois autres} toutes entierement differentes
entre elles; deux d’entre elles même peuvent être conduites de deux differentes
manieres chacune: ainsi je pourrois soutenir que je peus le demontrer de
\textls[100]{sept} manieres differentes. Les autres demonstrations que je prefererois pour
l’elegance aux deux données dans mon ouvrage seront publiées aussitot que
j’y trouverai l’occasion. A Propos, dans la premiere demonstration qui se
trouve dans la IV. section il s’est glissé une faute legere que je n’ai aperçue,
qu’apres que je ne pouvois plus l’indiquer. Il faut donc faire la correction
suivante.

p. 146 (cas (4)) l. 21 lisés comme il suit: »Facile vero perspicitur, ex
ista aequatione deduci posse haec $a'pRh\ldots(\alpha)$, $\pm ahRa'\ldots(\beta)$, $\pm ahRp\ldots(\gamma)$.
Ex $(\alpha)$ sequitur, perinde ut in (2), $h$ vel utriusque $a'$, $p$ vel neutrius residuum
esse. Sed casus prior ideo est impossibilis, quod ex $hRa'$ et $(\beta)$ sequeretur
$aRa'$ contra hypoth. Quamobrem necessario est $hNp$ adeoque, per $(\gamma)$, $aNp$.
Q. E. D.«[**)]
\par\vfill\noindent\rule{.25\linewidth}{.4pt}\par\smallskip
\begingroup\footnotesize
\noindent[*) Werke I, S. 87—94.]\par
\noindent[**) Werke I, S. 108 ist die Stelle nach der Anmerkung berichtigt, die \textsc{Gauss} dem art. 2 der Abhandlung \textit{Theorematis arithmetici demonstratio nova}, Comm. Soc. Reg. scient. Gotting. 16, 1808; Comm. classis mathem. ad annos 1801—1805, S. 70 beigefügt hat; vergl. Werke I, S. 477.]\par
\endgroup
\noindent X$_1$.\hfill10\par\clearpage
\SourceLeaf{081}\ArithmeticHead{74}{L}{ARITHMETIK.\quad NACHLASS.}
\par\indent Au reste à la page 144 il se trouve une faute d’impression non-indiquée,
savoir art. 139 ligne 3 au lieu de $\pm aNp$ il faut lire $\pm aRp$[*)].

J’aurois repondu plus tot à Votre lettre, mais la decouverte d’une nouvelle
planete par M. \textsc{Olbers} m’a un peu distrait. Par le premier essai que
j’ai fait sur son orbite, je trouve son mouvement considerablement plus vite
que celui de Cérès, Pallas et Junon, savoir $978''$ par jour. L’inclinaison de
l’orbite de $7^0\,6'$. L’excentricite 0,1. Cette planete a beaucoup plus de clarté
que Cérès, Pallas et Junon, et j’espere la trouver parmi les observations de
l’histoire celeste, peut être même parmi celles de \textsc{Flamsteed}. Je viens
d’achever un ouvrage etendu sur les methodes, qui me sont propres, à determiner
les orbites des planetes. Mais quoique je l’aie ecrit en allemand, je
trouve beaucoup de difficulté d’y engager un libraire.

La guerre a suspendu tout commerce, plusieurs de nos plus grands libraires
l’ont refusé. Je suis à present à traiter avec un autre qui se montre un peu
plus courageux. S’il trouvera son conte à cette entreprise, peut être il sera
encouragé par-là à risquer la publication d’un second volume de mes \textit{disquisitiones}.

Continuez, Mademoiselle, de me favoriser de Votre amitié et de Votre
correspondance, qui font mon orgueil, et soïes persuadée, que je suis et serai
toujours avec la plus haute estime
\begin{flushright}
Votre plus sincère admirateur\\
\textsc{Ch. Fr. Gauss.}
\end{flushright}
\noindent\begin{tabular}{c}Bronsvic ce 30 April 1807\\jour de ma naissance.\end{tabular}
\ClosingRule
